% Options for packages loaded elsewhere
\PassOptionsToPackage{unicode}{hyperref}
\PassOptionsToPackage{hyphens}{url}
%
\documentclass[
]{article}
\usepackage{amsmath,amssymb}
\usepackage{lmodern}
\usepackage{iftex}
\ifPDFTeX
  \usepackage[T1]{fontenc}
  \usepackage[utf8]{inputenc}
  \usepackage{textcomp} % provide euro and other symbols
\else % if luatex or xetex
  \usepackage{unicode-math}
  \defaultfontfeatures{Scale=MatchLowercase}
  \defaultfontfeatures[\rmfamily]{Ligatures=TeX,Scale=1}
\fi
% Use upquote if available, for straight quotes in verbatim environments
\IfFileExists{upquote.sty}{\usepackage{upquote}}{}
\IfFileExists{microtype.sty}{% use microtype if available
  \usepackage[]{microtype}
  \UseMicrotypeSet[protrusion]{basicmath} % disable protrusion for tt fonts
}{}
\makeatletter
\@ifundefined{KOMAClassName}{% if non-KOMA class
  \IfFileExists{parskip.sty}{%
    \usepackage{parskip}
  }{% else
    \setlength{\parindent}{0pt}
    \setlength{\parskip}{6pt plus 2pt minus 1pt}}
}{% if KOMA class
  \KOMAoptions{parskip=half}}
\makeatother
\usepackage{xcolor}
\usepackage[margin=1in]{geometry}
\usepackage{color}
\usepackage{fancyvrb}
\newcommand{\VerbBar}{|}
\newcommand{\VERB}{\Verb[commandchars=\\\{\}]}
\DefineVerbatimEnvironment{Highlighting}{Verbatim}{commandchars=\\\{\}}
% Add ',fontsize=\small' for more characters per line
\usepackage{framed}
\definecolor{shadecolor}{RGB}{248,248,248}
\newenvironment{Shaded}{\begin{snugshade}}{\end{snugshade}}
\newcommand{\AlertTok}[1]{\textcolor[rgb]{0.94,0.16,0.16}{#1}}
\newcommand{\AnnotationTok}[1]{\textcolor[rgb]{0.56,0.35,0.01}{\textbf{\textit{#1}}}}
\newcommand{\AttributeTok}[1]{\textcolor[rgb]{0.77,0.63,0.00}{#1}}
\newcommand{\BaseNTok}[1]{\textcolor[rgb]{0.00,0.00,0.81}{#1}}
\newcommand{\BuiltInTok}[1]{#1}
\newcommand{\CharTok}[1]{\textcolor[rgb]{0.31,0.60,0.02}{#1}}
\newcommand{\CommentTok}[1]{\textcolor[rgb]{0.56,0.35,0.01}{\textit{#1}}}
\newcommand{\CommentVarTok}[1]{\textcolor[rgb]{0.56,0.35,0.01}{\textbf{\textit{#1}}}}
\newcommand{\ConstantTok}[1]{\textcolor[rgb]{0.00,0.00,0.00}{#1}}
\newcommand{\ControlFlowTok}[1]{\textcolor[rgb]{0.13,0.29,0.53}{\textbf{#1}}}
\newcommand{\DataTypeTok}[1]{\textcolor[rgb]{0.13,0.29,0.53}{#1}}
\newcommand{\DecValTok}[1]{\textcolor[rgb]{0.00,0.00,0.81}{#1}}
\newcommand{\DocumentationTok}[1]{\textcolor[rgb]{0.56,0.35,0.01}{\textbf{\textit{#1}}}}
\newcommand{\ErrorTok}[1]{\textcolor[rgb]{0.64,0.00,0.00}{\textbf{#1}}}
\newcommand{\ExtensionTok}[1]{#1}
\newcommand{\FloatTok}[1]{\textcolor[rgb]{0.00,0.00,0.81}{#1}}
\newcommand{\FunctionTok}[1]{\textcolor[rgb]{0.00,0.00,0.00}{#1}}
\newcommand{\ImportTok}[1]{#1}
\newcommand{\InformationTok}[1]{\textcolor[rgb]{0.56,0.35,0.01}{\textbf{\textit{#1}}}}
\newcommand{\KeywordTok}[1]{\textcolor[rgb]{0.13,0.29,0.53}{\textbf{#1}}}
\newcommand{\NormalTok}[1]{#1}
\newcommand{\OperatorTok}[1]{\textcolor[rgb]{0.81,0.36,0.00}{\textbf{#1}}}
\newcommand{\OtherTok}[1]{\textcolor[rgb]{0.56,0.35,0.01}{#1}}
\newcommand{\PreprocessorTok}[1]{\textcolor[rgb]{0.56,0.35,0.01}{\textit{#1}}}
\newcommand{\RegionMarkerTok}[1]{#1}
\newcommand{\SpecialCharTok}[1]{\textcolor[rgb]{0.00,0.00,0.00}{#1}}
\newcommand{\SpecialStringTok}[1]{\textcolor[rgb]{0.31,0.60,0.02}{#1}}
\newcommand{\StringTok}[1]{\textcolor[rgb]{0.31,0.60,0.02}{#1}}
\newcommand{\VariableTok}[1]{\textcolor[rgb]{0.00,0.00,0.00}{#1}}
\newcommand{\VerbatimStringTok}[1]{\textcolor[rgb]{0.31,0.60,0.02}{#1}}
\newcommand{\WarningTok}[1]{\textcolor[rgb]{0.56,0.35,0.01}{\textbf{\textit{#1}}}}
\usepackage{graphicx}
\makeatletter
\def\maxwidth{\ifdim\Gin@nat@width>\linewidth\linewidth\else\Gin@nat@width\fi}
\def\maxheight{\ifdim\Gin@nat@height>\textheight\textheight\else\Gin@nat@height\fi}
\makeatother
% Scale images if necessary, so that they will not overflow the page
% margins by default, and it is still possible to overwrite the defaults
% using explicit options in \includegraphics[width, height, ...]{}
\setkeys{Gin}{width=\maxwidth,height=\maxheight,keepaspectratio}
% Set default figure placement to htbp
\makeatletter
\def\fps@figure{htbp}
\makeatother
\setlength{\emergencystretch}{3em} % prevent overfull lines
\providecommand{\tightlist}{%
  \setlength{\itemsep}{0pt}\setlength{\parskip}{0pt}}
\setcounter{secnumdepth}{-\maxdimen} % remove section numbering
\ifLuaTeX
  \usepackage{selnolig}  % disable illegal ligatures
\fi
\IfFileExists{bookmark.sty}{\usepackage{bookmark}}{\usepackage{hyperref}}
\IfFileExists{xurl.sty}{\usepackage{xurl}}{} % add URL line breaks if available
\urlstyle{same} % disable monospaced font for URLs
\hypersetup{
  pdftitle={211-DiscussionI},
  pdfauthor={Yu-Shiuan (Lily) Huang},
  hidelinks,
  pdfcreator={LaTeX via pandoc}}

\title{211-DiscussionI}
\author{Yu-Shiuan (Lily) Huang}
\date{2023-10-03}

\begin{document}
\maketitle

\hypertarget{exploring-a-new-dataset-using-r-part-i}{%
\paragraph{Exploring a New Dataset using R (Part
I)}\label{exploring-a-new-dataset-using-r-part-i}}

Note: All these R codes are prepared by Prof.~Ryan Hubert.

\hypertarget{a.-load-r-packages}{%
\subparagraph{a. Load R packages}\label{a.-load-r-packages}}

There are many useful packages in R that can be installed freely. If it
is the first time you are going to use a package, here are the steps:

\begin{itemize}
\tightlist
\item
  Install it
\item
  library it (which means telling R you are going to use it)
\end{itemize}

Once you have installed the package, you don't have to install it every
time you open R studio. However, you still have to load it through the
\texttt{library} function every time.

\begin{Shaded}
\begin{Highlighting}[]
\CommentTok{\# Install packages}
\CommentTok{\# install.packages("tidyverse") \# manipulate data }
\CommentTok{\# install.packages("Matrix") \# calculate rank}
\DocumentationTok{\#\# these are the packages we will use for today }

\CommentTok{\# load packages}
\FunctionTok{library}\NormalTok{(tidyverse)}
\FunctionTok{library}\NormalTok{(Matrix) }
\end{Highlighting}
\end{Shaded}

The \texttt{tidyverse} is a powerful collection of R packages that are
actually data tools for transforming and visualizing data. All packages
of the \texttt{tidyverse} share an underlying philosophy and common
APls.

\includegraphics[width=0.6\textwidth,height=\textheight]{https://cdn.analyticsvidhya.com/wp-content/uploads/2019/05/ggplot_hive.jpg}

The core packages within the \texttt{tidyverse} are:

\begin{itemize}
\tightlist
\item
  \texttt{ggplot2}, which implements the grammar of graphics. You can
  use it to visualize your data.
\item
  \texttt{dplyr} is a grammar of data. You can use it to solve the most
  common data manipulation challenges.
\item
  \texttt{tidyr} helps you to create tidy data or data where each
  variable is in a column, each observation is a row end each value is a
  column, each observation is a row end each value is a cell.
\item
  \texttt{readr} is a fast and friendly way to read rectangular
\item
  \texttt{purrr} enhances R's functional programming (FP) toolkit by
  providing a complete and consistent set of tools for working with
  functions and vectors.
\item
  \texttt{tibble} is a modern re-imaginging of the data
\item
  \texttt{stringr} provides a cohesive set of functions designed to make
  working with strings as easy as possible
\item
  \texttt{forcats} provide a suite of useful tools that solve common
  problems with factors.
\end{itemize}

Highly recommend this
\href{https://www.rstudio.com/wp-content/uploads/2015/02/data-wrangling-cheatsheet.pdf}{data
wrangling cheat sheet} with \texttt{dplyr} and \texttt{tidyr} in R!

\hypertarget{b.-import-data}{%
\subparagraph{b. Import data}\label{b.-import-data}}

For today, we are using \texttt{USDC-DATASET-JOP.csv}. You can download
this csv file from Canvas or from
\href{https://dataverse.harvard.edu/file.xhtml?fileId=4306747\&version=1.0.}{here}.

\begin{Shaded}
\begin{Highlighting}[]
\CommentTok{\# Import data}
\NormalTok{df }\OtherTok{\textless{}{-}} \FunctionTok{read\_csv}\NormalTok{(}\StringTok{"/Users/yu{-}shiuanhuang/Desktop/method{-}sequence/data/USDC{-}DATASET{-}JOP.csv"}\NormalTok{) }\CommentTok{\# Here we assign the dataset as an object called “df”}

\DocumentationTok{\#\# Note: you create new objects by assigning them with the \textasciigrave{}\textless{}{-}\textasciigrave{} phrase. You can also use an equal sign \textasciigrave{}=\textasciigrave{}.}
\end{Highlighting}
\end{Shaded}

Here are some tips for how to find the file path of the data you are
going to use.

\begin{itemize}
\item
  First, where did you store your data? In which folder? In my case, I
  stored the data in the \texttt{data} folder within the
  \texttt{method-sequence} folder.
\item
  Second, use \texttt{getwd} function to get the grasp of how does a
  file path look like. \texttt{getwd} function returns an absolute file
  path representing the current working directory of the R process.
\end{itemize}

\begin{Shaded}
\begin{Highlighting}[]
  \FunctionTok{getwd}\NormalTok{()}
\end{Highlighting}
\end{Shaded}

\begin{verbatim}
## [1] "/Users/yu-shiuanhuang/Desktop/method-sequence/POL211"
\end{verbatim}

As you can see, my current working directory is
\texttt{"/Users/yu-shiuanhuang/Desktop/method-sequence/docs"}. Since I
know that I stored the data in the \texttt{data} folder within the
\texttt{method-sequence} folder, I can simply just change the file path
to
\texttt{"/Users/yu-shiuanhuang/Desktop/method-sequence/data/USDC-DATASET-JOP.csv"}
and put it in the \texttt{read\_csv} function so R knows where to read
the file you want.

If you import the data successfully, you should be able to see df
popping up in the environment section.

\includegraphics[width=0.6\textwidth,height=\textheight]{/Users/yu-shiuanhuang/Desktop/method-sequence/figures/environment.png}

Sometimes the dataset we are interested in may be stored in a different
form, such as \texttt{dta}, \texttt{spss}, \texttt{xlsx}, \texttt{rds},
\texttt{rdata}, etc. You can find the corresponded codes from
\href{https://www.datafiles.samhsa.gov/get-help/format-specific-issues/how-do-i-read-data-r}{here}.

\hypertarget{c.-basic-information-about-the-dataset}{%
\subparagraph{c.~Basic information about the
dataset}\label{c.-basic-information-about-the-dataset}}

What is the size of the dataset?

\begin{Shaded}
\begin{Highlighting}[]
\CommentTok{\# the number of row is:}
\FunctionTok{nrow}\NormalTok{(df)}
\end{Highlighting}
\end{Shaded}

\begin{verbatim}
## [1] 97725
\end{verbatim}

How many observations (rows) and variables (columns) we have in the
dataset?

\begin{Shaded}
\begin{Highlighting}[]
\CommentTok{\# the number of column is:}
\FunctionTok{ncol}\NormalTok{(df)}
\end{Highlighting}
\end{Shaded}

\begin{verbatim}
## [1] 54
\end{verbatim}

You can also use \texttt{dim} function to get the size of the dataset.

\begin{Shaded}
\begin{Highlighting}[]
\FunctionTok{dim}\NormalTok{(df) }
\end{Highlighting}
\end{Shaded}

\begin{verbatim}
## [1] 97725    54
\end{verbatim}

\begin{Shaded}
\begin{Highlighting}[]
\DocumentationTok{\#\# this function gives you the number of rows and columns at the same time}
\end{Highlighting}
\end{Shaded}

The size of this dataset is 97,725 x 54.

What does the dataset actually look like?

If you want to just quickly see the first six rows to get a feel for
what the dataset looks like:

\begin{Shaded}
\begin{Highlighting}[]
\FunctionTok{head}\NormalTok{(df)}
\end{Highlighting}
\end{Shaded}

\begin{verbatim}
## # A tibble: 6 x 54
##   file      OUTCOME jrepublican jpresident jid_anon jyears APPEAL judge_replaced
##   <chr>     <chr>         <dbl> <chr>         <dbl>  <dbl>  <dbl>          <dbl>
## 1 /dockets~ jud_de~           1 Reagan        90059     13      0              0
## 2 /dockets~ dism_o~           1 Reagan        90048      9      0              0
## 3 /dockets~ dism_p~           1 Reagan        90012     11      0              0
## 4 /dockets~ dism_o~           1 Bush41        90016      3      0              0
## 5 /dockets~ dism_v~           1 Reagan        90096     10      0              0
## 6 /dockets~ dism_v~           1 Reagan        90158     11      0              0
## # i 46 more variables: jsenior <dbl>, court <chr>, division <dbl>, YEAR <dbl>,
## #   QUARTER <chr>, nature_of_suit <dbl>, SECTION <chr>, ORIGIN <dbl>,
## #   JURIS <chr>, jury_demand <chr>, PROSE <dbl>, COUNTY <chr>, pla_count <dbl>,
## #   pla_count_prose <dbl>, pla_count_anonymous <dbl>, pla_count_IND <dbl>,
## #   pla_count_BUS <dbl>, pla_count_GOV <dbl>, pla_count_LAW <dbl>,
## #   pla_count_OTH <dbl>, pla_count_repeat <dbl>, pla_acount <dbl>,
## #   pla_acount_repeat <dbl>, def_count <dbl>, def_count_prose <dbl>, ...
\end{verbatim}

Note that a lot of the variables are not visible because R will try to
clean up output.

One convenient thing is that you can see the type of each
variable.\href{https://tibble.tidyverse.org/articles/types.html}{Here}
is the full list of all the possible types of variables (in tidyverse).

If you want to look at the whole dataset, use \texttt{View(df)} command
and it will open in a new tab in R.

\hypertarget{d.-exploring-the-variables}{%
\subparagraph{d.~Exploring the
variables}\label{d.-exploring-the-variables}}

Most people prefer to look at variables using their names.

How to ask R to tell us the name of each variables?

\begin{Shaded}
\begin{Highlighting}[]
\CommentTok{\# Name of each column}
\FunctionTok{colnames}\NormalTok{(df)}
\end{Highlighting}
\end{Shaded}

\begin{verbatim}
##  [1] "file"                "OUTCOME"             "jrepublican"        
##  [4] "jpresident"          "jid_anon"            "jyears"             
##  [7] "APPEAL"              "judge_replaced"      "jsenior"            
## [10] "court"               "division"            "YEAR"               
## [13] "QUARTER"             "nature_of_suit"      "SECTION"            
## [16] "ORIGIN"              "JURIS"               "jury_demand"        
## [19] "PROSE"               "COUNTY"              "pla_count"          
## [22] "pla_count_prose"     "pla_count_anonymous" "pla_count_IND"      
## [25] "pla_count_BUS"       "pla_count_GOV"       "pla_count_LAW"      
## [28] "pla_count_OTH"       "pla_count_repeat"    "pla_acount"         
## [31] "pla_acount_repeat"   "def_count"           "def_count_prose"    
## [34] "def_count_anonymous" "def_count_IND"       "def_count_BUS"      
## [37] "def_count_GOV"       "def_count_LAW"       "def_count_OTH"      
## [40] "def_count_repeat"    "def_acount"          "def_acount_repeat"  
## [43] "oth_count"           "oth_count_prose"     "oth_count_anonymous"
## [46] "oth_count_IND"       "oth_count_BUS"       "oth_count_GOV"      
## [49] "oth_count_LAW"       "oth_count_OTH"       "oth_count_repeat"   
## [52] "oth_acount"          "oth_acount_repeat"   "ifp_denial"
\end{verbatim}

How to call out one specific variable we are interested in from the
dataset? There are multiple ways to look at an arbitrary variable:

\begin{Shaded}
\begin{Highlighting}[]
\DocumentationTok{\#\# Use column number}
\NormalTok{df[, }\DecValTok{2}\NormalTok{] }
\NormalTok{df[[}\DecValTok{2}\NormalTok{]]}

\DocumentationTok{\#\# Use variable name}
\NormalTok{df[, }\StringTok{"OUTCOME"}\NormalTok{]}
\NormalTok{df[[}\StringTok{"OUTCOME"}\NormalTok{]]}
\NormalTok{df}\SpecialCharTok{$}\NormalTok{OUTCOME }\DocumentationTok{\#\# This is the most common way to do it}
\end{Highlighting}
\end{Shaded}

\begin{Shaded}
\begin{Highlighting}[]
\FunctionTok{head}\NormalTok{(df}\SpecialCharTok{$}\NormalTok{OUTCOME) }\DocumentationTok{\#\# To see the first six rows of this variable}
\end{Highlighting}
\end{Shaded}

\begin{verbatim}
## [1] "jud_default_plaintiff" "dism_other"            "dism_pros"            
## [4] "dism_other"            "dism_vol"              "dism_vol"
\end{verbatim}

What kind of variable is the \texttt{OUTCOME} variable? It looks like a
categorical variable!

The next question would be what are the categories (called ``levels'' in
R) of this variable?

\begin{Shaded}
\begin{Highlighting}[]
\CommentTok{\# This asks R to give us a list of all the values used in this variable.}
\FunctionTok{unique}\NormalTok{(df}\SpecialCharTok{$}\NormalTok{OUTCOME)}
\end{Highlighting}
\end{Shaded}

\begin{verbatim}
##  [1] "jud_default_plaintiff"  "dism_other"             "dism_pros"             
##  [4] "dism_vol"               "settlement"             "jud_motion_defendant"  
##  [7] "dism_juris"             "remand"                 "jud_misc_NA"           
## [10] "jud_misc_both"          "jud_motion_plaintiff"   "jud_misc_defendant"    
## [13] "jud_jury_defendant"     "transfer"               "jud_bench_defendant"   
## [16] "jud_misc_plaintiff"     "jud_consent_plaintiff"  "jud_jury_plaintiff"    
## [19] "jud_default_both"       "jud_motion_both"        "jud_default_defendant" 
## [22] "jud_bench_plaintiff"    "jud_jury_both"          "jud_consent_both"      
## [25] "jud_bench_both"         "jud_directed_defendant" "jud_consent_defendant" 
## [28] "jud_motion_NA"          "jud_jury_NA"            "jud_bench_NA"          
## [31] "jud_consent_NA"         "jud_directed_plaintiff" "jud_default_NA"        
## [34] "jud_directed_NA"
\end{verbatim}

There are 34 categories in this variable.

Here's an issue: when R loaded the dataset, it didn't code this variable
as a categorical variable. Since the \texttt{OUTCOME} variable is coded
as strings/texts in the dataset, in R, the type of this variable is
character.

\begin{Shaded}
\begin{Highlighting}[]
\CommentTok{\# Look at what this says about the variable type:}
\FunctionTok{type\_sum}\NormalTok{(df}\SpecialCharTok{$}\NormalTok{OUTCOME)}
\end{Highlighting}
\end{Shaded}

\begin{verbatim}
## [1] "chr"
\end{verbatim}

\begin{Shaded}
\begin{Highlighting}[]
\CommentTok{\# When you summarize it in R, R doesn’t provide you with meaningful information}
\FunctionTok{summary}\NormalTok{(df}\SpecialCharTok{$}\NormalTok{OUTCOME)}
\end{Highlighting}
\end{Shaded}

\begin{verbatim}
##    Length     Class      Mode 
##     97725 character character
\end{verbatim}

Good news is that we can transform the type of a variable in R.

\begin{Shaded}
\begin{Highlighting}[]
\CommentTok{\# lets\textquotesingle{}s change the type to a categorical variable, known as a "factor" variable in R:}
\NormalTok{df}\SpecialCharTok{$}\NormalTok{OUTCOME }\OtherTok{\textless{}{-}} \FunctionTok{as.factor}\NormalTok{(df}\SpecialCharTok{$}\NormalTok{OUTCOME) }
\DocumentationTok{\#\# In order to override the original variable, we have to assign the new version of the variable to the original one.}

\CommentTok{\# After factorizing, R can summarize in a more meaningful way: what is the number of each category?}
\FunctionTok{summary}\NormalTok{(df}\SpecialCharTok{$}\NormalTok{OUTCOME)}
\end{Highlighting}
\end{Shaded}

\begin{verbatim}
##             dism_juris             dism_other              dism_pros 
##                   1321                  16198                   3123 
##               dism_vol         jud_bench_both    jud_bench_defendant 
##                  17805                     10                    192 
##           jud_bench_NA    jud_bench_plaintiff       jud_consent_both 
##                     22                    110                    199 
##  jud_consent_defendant         jud_consent_NA  jud_consent_plaintiff 
##                     48                     93                    391 
##       jud_default_both  jud_default_defendant         jud_default_NA 
##                      5                     69                      7 
##  jud_default_plaintiff jud_directed_defendant        jud_directed_NA 
##                    228                     34                      1 
## jud_directed_plaintiff          jud_jury_both     jud_jury_defendant 
##                      6                     73                   1260 
##            jud_jury_NA     jud_jury_plaintiff          jud_misc_both 
##                    121                    528                    399 
##     jud_misc_defendant            jud_misc_NA     jud_misc_plaintiff 
##                   1355                   2239                    413 
##        jud_motion_both   jud_motion_defendant          jud_motion_NA 
##                    428                  10510                    821 
##   jud_motion_plaintiff                 remand             settlement 
##                    629                   4546                  32627 
##               transfer 
##                   1914
\end{verbatim}

\hypertarget{e.-exploring-the-observations}{%
\subparagraph{e. Exploring the
observations}\label{e.-exploring-the-observations}}

Each row (observation) in this dataset is a court case. We can also use
R to look at an arbitrary one.

\begin{Shaded}
\begin{Highlighting}[]
\NormalTok{df[}\DecValTok{5}\NormalTok{,] }\CommentTok{\# This means that we want to call out the fifth row of the data.}
\end{Highlighting}
\end{Shaded}

\begin{verbatim}
## # A tibble: 1 x 54
##   file      OUTCOME jrepublican jpresident jid_anon jyears APPEAL judge_replaced
##   <chr>     <fct>         <dbl> <chr>         <dbl>  <dbl>  <dbl>          <dbl>
## 1 /dockets~ dism_v~           1 Reagan        90096     10      0              0
## # i 46 more variables: jsenior <dbl>, court <chr>, division <dbl>, YEAR <dbl>,
## #   QUARTER <chr>, nature_of_suit <dbl>, SECTION <chr>, ORIGIN <dbl>,
## #   JURIS <chr>, jury_demand <chr>, PROSE <dbl>, COUNTY <chr>, pla_count <dbl>,
## #   pla_count_prose <dbl>, pla_count_anonymous <dbl>, pla_count_IND <dbl>,
## #   pla_count_BUS <dbl>, pla_count_GOV <dbl>, pla_count_LAW <dbl>,
## #   pla_count_OTH <dbl>, pla_count_repeat <dbl>, pla_acount <dbl>,
## #   pla_acount_repeat <dbl>, def_count <dbl>, def_count_prose <dbl>, ...
\end{verbatim}

\begin{Shaded}
\begin{Highlighting}[]
\CommentTok{\# If you want to see all the variables for that observation, here are some options:}
\FunctionTok{t}\NormalTok{(df[}\DecValTok{5}\NormalTok{,]) }\DocumentationTok{\#\# downside: this treats all the values like strings!}
\end{Highlighting}
\end{Shaded}

\begin{verbatim}
##                     [,1]                                  
## file                "/dockets-cacd/1995/199501-00068.html"
## OUTCOME             "dism_vol"                            
## jrepublican         "1"                                   
## jpresident          "Reagan"                              
## jid_anon            "90096"                               
## jyears              "10"                                  
## APPEAL              "0"                                   
## judge_replaced      "0"                                   
## jsenior             "0"                                   
## court               "cacd"                                
## division            "2"                                   
## YEAR                "1995"                                
## QUARTER             "1995Q1"                              
## nature_of_suit      "442"                                 
## SECTION             "1983"                                
## ORIGIN              "1"                                   
## JURIS               "federal_question"                    
## jury_demand         "Plaintiff"                           
## PROSE               NA                                    
## COUNTY              "06037"                               
## pla_count           "1"                                   
## pla_count_prose     "0"                                   
## pla_count_anonymous "0"                                   
## pla_count_IND       "0.7538"                              
## pla_count_BUS       "0.1453"                              
## pla_count_GOV       "0.0485"                              
## pla_count_LAW       "0.0315"                              
## pla_count_OTH       "0.021"                               
## pla_count_repeat    "0"                                   
## pla_acount          "1"                                   
## pla_acount_repeat   "1"                                   
## def_count           "4"                                   
## def_count_prose     "0"                                   
## def_count_anonymous "1"                                   
## def_count_IND       "2.7986"                              
## def_count_BUS       "0.7653"                              
## def_count_GOV       "0.2055"                              
## def_count_LAW       "0.1402"                              
## def_count_OTH       "0.0903"                              
## def_count_repeat    "1"                                   
## def_acount          "0"                                   
## def_acount_repeat   "0"                                   
## oth_count           "0"                                   
## oth_count_prose     "0"                                   
## oth_count_anonymous "0"                                   
## oth_count_IND       "0"                                   
## oth_count_BUS       "0"                                   
## oth_count_GOV       "0"                                   
## oth_count_LAW       "0"                                   
## oth_count_OTH       "0"                                   
## oth_count_repeat    "0"                                   
## oth_acount          "0"                                   
## oth_acount_repeat   "0"                                   
## ifp_denial          "0"
\end{verbatim}

\begin{Shaded}
\begin{Highlighting}[]
\FunctionTok{as.data.frame}\NormalTok{(df[}\DecValTok{5}\NormalTok{,]) }\DocumentationTok{\#\# downside: sort of ugly and hard to see}
\end{Highlighting}
\end{Shaded}

\begin{verbatim}
##                                   file  OUTCOME jrepublican jpresident jid_anon
## 1 /dockets-cacd/1995/199501-00068.html dism_vol           1     Reagan    90096
##   jyears APPEAL judge_replaced jsenior court division YEAR QUARTER
## 1     10      0              0       0  cacd        2 1995  1995Q1
##   nature_of_suit SECTION ORIGIN            JURIS jury_demand PROSE COUNTY
## 1            442    1983      1 federal_question   Plaintiff    NA  06037
##   pla_count pla_count_prose pla_count_anonymous pla_count_IND pla_count_BUS
## 1         1               0                   0        0.7538        0.1453
##   pla_count_GOV pla_count_LAW pla_count_OTH pla_count_repeat pla_acount
## 1        0.0485        0.0315         0.021                0          1
##   pla_acount_repeat def_count def_count_prose def_count_anonymous def_count_IND
## 1                 1         4               0                   1        2.7986
##   def_count_BUS def_count_GOV def_count_LAW def_count_OTH def_count_repeat
## 1        0.7653        0.2055        0.1402        0.0903                1
##   def_acount def_acount_repeat oth_count oth_count_prose oth_count_anonymous
## 1          0                 0         0               0                   0
##   oth_count_IND oth_count_BUS oth_count_GOV oth_count_LAW oth_count_OTH
## 1             0             0             0             0             0
##   oth_count_repeat oth_acount oth_acount_repeat ifp_denial
## 1                0          0                 0          0
\end{verbatim}

\begin{Shaded}
\begin{Highlighting}[]
\DocumentationTok{\#\# To draw out an element from a data, you have to specify a row number and a column number.}
\NormalTok{df[}\DecValTok{5}\NormalTok{, }\DecValTok{1}\NormalTok{]}
\end{Highlighting}
\end{Shaded}

\begin{verbatim}
## # A tibble: 1 x 1
##   file                                
##   <chr>                               
## 1 /dockets-cacd/1995/199501-00068.html
\end{verbatim}

Usually, you only need to see some of the variables. For example:
suppose we want to know which judge heard this case
(\texttt{jid\_anon}), what the case outcome was (\texttt{OUTCOME}) and
the party of the judge's appointing president (\texttt{jrepublican}).

\begin{Shaded}
\begin{Highlighting}[]
\NormalTok{df[}\DecValTok{5}\NormalTok{, }\FunctionTok{c}\NormalTok{(}\StringTok{"jid\_anon"}\NormalTok{, }\StringTok{"OUTCOME"}\NormalTok{, }\StringTok{"jrepublican"}\NormalTok{)]}
\end{Highlighting}
\end{Shaded}

\begin{verbatim}
## # A tibble: 1 x 3
##   jid_anon OUTCOME  jrepublican
##      <dbl> <fct>          <dbl>
## 1    90096 dism_vol           1
\end{verbatim}

A more neat way to write the code by using \texttt{tidyverse} package.

\begin{Shaded}
\begin{Highlighting}[]
\CommentTok{\# using pipes to do the same thing; use select function from tidyverse to pick variables }
\NormalTok{df[}\DecValTok{5}\NormalTok{, ] }\SpecialCharTok{\%\textgreater{}\%} \FunctionTok{select}\NormalTok{(jid\_anon, OUTCOME, jrepublican)}
\end{Highlighting}
\end{Shaded}

\begin{verbatim}
## # A tibble: 1 x 3
##   jid_anon OUTCOME  jrepublican
##      <dbl> <fct>          <dbl>
## 1    90096 dism_vol           1
\end{verbatim}

\texttt{\%\textgreater{}\%} is called ``pipes.'' Pipes take the output
from one function and feed it to the first argument of the next
function.

We can also find a row by ``filtering'' the data based on some criteria.

What if I want to find a case by its ID number? Eg,
``/dockets-cacd/2007/200738-00195.html''

\begin{Shaded}
\begin{Highlighting}[]
\CommentTok{\# use filter function from tidyverse to set the condition }
\NormalTok{df }\SpecialCharTok{\%\textgreater{}\%} \FunctionTok{filter}\NormalTok{(file }\SpecialCharTok{==} \StringTok{"/dockets{-}cacd/2007/200738{-}00195.html"}\NormalTok{)}
\end{Highlighting}
\end{Shaded}

\begin{verbatim}
## # A tibble: 1 x 54
##   file      OUTCOME jrepublican jpresident jid_anon jyears APPEAL judge_replaced
##   <chr>     <fct>         <dbl> <chr>         <dbl>  <dbl>  <dbl>          <dbl>
## 1 /dockets~ settle~           1 Bush43        90104      0      0              0
## # i 46 more variables: jsenior <dbl>, court <chr>, division <dbl>, YEAR <dbl>,
## #   QUARTER <chr>, nature_of_suit <dbl>, SECTION <chr>, ORIGIN <dbl>,
## #   JURIS <chr>, jury_demand <chr>, PROSE <dbl>, COUNTY <chr>, pla_count <dbl>,
## #   pla_count_prose <dbl>, pla_count_anonymous <dbl>, pla_count_IND <dbl>,
## #   pla_count_BUS <dbl>, pla_count_GOV <dbl>, pla_count_LAW <dbl>,
## #   pla_count_OTH <dbl>, pla_count_repeat <dbl>, pla_acount <dbl>,
## #   pla_acount_repeat <dbl>, def_count <dbl>, def_count_prose <dbl>, ...
\end{verbatim}

What if I want to find all cases heard by judges appointed by Bush 43?

\begin{Shaded}
\begin{Highlighting}[]
\NormalTok{df }\SpecialCharTok{\%\textgreater{}\%} \FunctionTok{filter}\NormalTok{(jpresident }\SpecialCharTok{==} \StringTok{"Bush43"}\NormalTok{)}
\end{Highlighting}
\end{Shaded}

\begin{verbatim}
## # A tibble: 17,622 x 54
##    file     OUTCOME jrepublican jpresident jid_anon jyears APPEAL judge_replaced
##    <chr>    <fct>         <dbl> <chr>         <dbl>  <dbl>  <dbl>          <dbl>
##  1 /docket~ jud_mo~           1 Bush43        90131      1      0              0
##  2 /docket~ dism_v~           1 Bush43        90131      0      0              0
##  3 /docket~ dism_p~           1 Bush43        90131      0      0              0
##  4 /docket~ dism_j~           1 Bush43        90131      0      0              0
##  5 /docket~ dism_o~           1 Bush43        90159      0      0              0
##  6 /docket~ dism_p~           1 Bush43        90159      0      0              0
##  7 /docket~ dism_o~           1 Bush43        90159      0      0              0
##  8 /docket~ dism_o~           1 Bush43        90159      0      0              0
##  9 /docket~ dism_o~           1 Bush43        90131      0      1              0
## 10 /docket~ dism_o~           1 Bush43        90159      0      1              1
## # i 17,612 more rows
## # i 46 more variables: jsenior <dbl>, court <chr>, division <dbl>, YEAR <dbl>,
## #   QUARTER <chr>, nature_of_suit <dbl>, SECTION <chr>, ORIGIN <dbl>,
## #   JURIS <chr>, jury_demand <chr>, PROSE <dbl>, COUNTY <chr>, pla_count <dbl>,
## #   pla_count_prose <dbl>, pla_count_anonymous <dbl>, pla_count_IND <dbl>,
## #   pla_count_BUS <dbl>, pla_count_GOV <dbl>, pla_count_LAW <dbl>,
## #   pla_count_OTH <dbl>, pla_count_repeat <dbl>, pla_acount <dbl>, ...
\end{verbatim}

The reason why we need to add quotation mark when setting the condition
in the \texttt{filter} function is that both of the \texttt{file} and
\texttt{jpresident} variables are coded in character type (categorical)
in R. If we filter our cases based on whether the judge of the case is
assigned by a Republican president (when \texttt{jrepublican} == 1), we
don't have to add quotation mark in the condition as this variable is
coded as the type of double (dbl) in R, which is a class of numeric.

\begin{Shaded}
\begin{Highlighting}[]
\NormalTok{df }\SpecialCharTok{\%\textgreater{}\%} \FunctionTok{filter}\NormalTok{(jrepublican }\SpecialCharTok{==} \DecValTok{1}\NormalTok{)}
\end{Highlighting}
\end{Shaded}

\begin{verbatim}
## # A tibble: 42,679 x 54
##    file     OUTCOME jrepublican jpresident jid_anon jyears APPEAL judge_replaced
##    <chr>    <fct>         <dbl> <chr>         <dbl>  <dbl>  <dbl>          <dbl>
##  1 /docket~ jud_de~           1 Reagan        90059     13      0              0
##  2 /docket~ dism_o~           1 Reagan        90048      9      0              0
##  3 /docket~ dism_p~           1 Reagan        90012     11      0              0
##  4 /docket~ dism_o~           1 Bush41        90016      3      0              0
##  5 /docket~ dism_v~           1 Reagan        90096     10      0              0
##  6 /docket~ dism_v~           1 Reagan        90158     11      0              0
##  7 /docket~ settle~           1 Reagan        90158     11      0              0
##  8 /docket~ settle~           1 Ford          90218     19      0              0
##  9 /docket~ jud_mo~           1 Nixon         90148     24      0              0
## 10 /docket~ dism_j~           1 Reagan        90005     13      0              0
## # i 42,669 more rows
## # i 46 more variables: jsenior <dbl>, court <chr>, division <dbl>, YEAR <dbl>,
## #   QUARTER <chr>, nature_of_suit <dbl>, SECTION <chr>, ORIGIN <dbl>,
## #   JURIS <chr>, jury_demand <chr>, PROSE <dbl>, COUNTY <chr>, pla_count <dbl>,
## #   pla_count_prose <dbl>, pla_count_anonymous <dbl>, pla_count_IND <dbl>,
## #   pla_count_BUS <dbl>, pla_count_GOV <dbl>, pla_count_LAW <dbl>,
## #   pla_count_OTH <dbl>, pla_count_repeat <dbl>, pla_acount <dbl>, ...
\end{verbatim}

\hypertarget{f.-some-measurement-issues}{%
\subparagraph{f.~Some measurement
issues}\label{f.-some-measurement-issues}}

What kinds of variable do we have?

There are slight differences between the way a computer
stores/understands data and the way we conceptualize it. In this
dataset, we have some variables on cardinal scales.

For example: How many plaintiffs are named in each case? (Why is this on
a cardinal scale?)

\begin{Shaded}
\begin{Highlighting}[]
\FunctionTok{head}\NormalTok{(df}\SpecialCharTok{$}\NormalTok{pla\_count)}
\end{Highlighting}
\end{Shaded}

\begin{verbatim}
## [1] 1 1 1 1 1 1
\end{verbatim}

What is the data type? It's ``dbl'' which means it's a number that need
not be an integer (can be fractional number, 3.51)

\begin{Shaded}
\begin{Highlighting}[]
\FunctionTok{type\_sum}\NormalTok{(df}\SpecialCharTok{$}\NormalTok{pla\_count)}
\end{Highlighting}
\end{Shaded}

\begin{verbatim}
## [1] "dbl"
\end{verbatim}

Note: this variable is actually comprised of integers, we could convert
it if we want, but this doesn't really matter. (Why? 11.0 = 11)

\begin{Shaded}
\begin{Highlighting}[]
\NormalTok{df}\SpecialCharTok{$}\NormalTok{pla\_count }\OtherTok{\textless{}{-}} \FunctionTok{as.integer}\NormalTok{(df}\SpecialCharTok{$}\NormalTok{pla\_count)}
\FunctionTok{type\_sum}\NormalTok{(df}\SpecialCharTok{$}\NormalTok{pla\_count)}
\end{Highlighting}
\end{Shaded}

\begin{verbatim}
## [1] "int"
\end{verbatim}

We also have some variables that are ordinal.

For example, in what years were cases filed? (Why is this on an ordinal
scale?)

\begin{Shaded}
\begin{Highlighting}[]
\FunctionTok{head}\NormalTok{(df}\SpecialCharTok{$}\NormalTok{YEAR)}
\end{Highlighting}
\end{Shaded}

\begin{verbatim}
## [1] 1995 1995 1995 1995 1995 1995
\end{verbatim}

What is the data type? It's also ``dbl''.

\begin{Shaded}
\begin{Highlighting}[]
\FunctionTok{type\_sum}\NormalTok{(df}\SpecialCharTok{$}\NormalTok{YEAR)}
\end{Highlighting}
\end{Shaded}

\begin{verbatim}
## [1] "dbl"
\end{verbatim}

Acknowledge, this is bad! We do not want R to treat this as a cardinal
variable.

For example, it (incorrectly) allows us to add them:

\begin{Shaded}
\begin{Highlighting}[]
\NormalTok{df}\SpecialCharTok{$}\NormalTok{YEAR[}\DecValTok{1}\NormalTok{] }\SpecialCharTok{+}\NormalTok{ df}\SpecialCharTok{$}\NormalTok{YEAR[}\DecValTok{2}\NormalTok{]}
\end{Highlighting}
\end{Shaded}

\begin{verbatim}
## [1] 3990
\end{verbatim}

\begin{Shaded}
\begin{Highlighting}[]
\DocumentationTok{\#\# Notice here the change in indexing! Remember that when you call out one variable from the data, it is a vector. To draw out an element from a vector, you only have to specify one number inside [ ].}
\end{Highlighting}
\end{Shaded}

So, we should change it to a factor variable so that we don't make this
mistake.

\begin{Shaded}
\begin{Highlighting}[]
\NormalTok{df}\SpecialCharTok{$}\NormalTok{YEAR }\OtherTok{\textless{}{-}} \FunctionTok{as.factor}\NormalTok{(df}\SpecialCharTok{$}\NormalTok{YEAR)}
\NormalTok{df}\SpecialCharTok{$}\NormalTok{YEAR[}\DecValTok{1}\NormalTok{] }\SpecialCharTok{+}\NormalTok{ df}\SpecialCharTok{$}\NormalTok{YEAR[}\DecValTok{2}\NormalTok{] }
\end{Highlighting}
\end{Shaded}

\begin{verbatim}
## [1] NA
\end{verbatim}

We got an error when we tried to do something silly, which is good!

In the analysis for Ryan's paper, they had to convert the
\texttt{OUTCOME} variable to dummies. How can we do it in R?

Recall that there are 34 unique categories in the \texttt{OUTCOME}
variable.

\begin{Shaded}
\begin{Highlighting}[]
\FunctionTok{unique}\NormalTok{(df}\SpecialCharTok{$}\NormalTok{OUTCOME) }
\end{Highlighting}
\end{Shaded}

\begin{verbatim}
##  [1] jud_default_plaintiff  dism_other             dism_pros             
##  [4] dism_vol               settlement             jud_motion_defendant  
##  [7] dism_juris             remand                 jud_misc_NA           
## [10] jud_misc_both          jud_motion_plaintiff   jud_misc_defendant    
## [13] jud_jury_defendant     transfer               jud_bench_defendant   
## [16] jud_misc_plaintiff     jud_consent_plaintiff  jud_jury_plaintiff    
## [19] jud_default_both       jud_motion_both        jud_default_defendant 
## [22] jud_bench_plaintiff    jud_jury_both          jud_consent_both      
## [25] jud_bench_both         jud_directed_defendant jud_consent_defendant 
## [28] jud_motion_NA          jud_jury_NA            jud_bench_NA          
## [31] jud_consent_NA         jud_directed_plaintiff jud_default_NA        
## [34] jud_directed_NA       
## 34 Levels: dism_juris dism_other dism_pros dism_vol ... transfer
\end{verbatim}

When we want to make each category of this variable into dummies, it
means that we need to create 34 new dummy variables (34 new columns for
the dataset). For example, the new dummy variable \texttt{dism\_other}
will be coded as 1 if the \texttt{OUTCOME} variable is coded as
dism\_other. Let's use \texttt{ifesle} function to create this dummy!

\begin{Shaded}
\begin{Highlighting}[]
\NormalTok{df}\SpecialCharTok{$}\NormalTok{dism\_other }\OtherTok{\textless{}{-}} \FunctionTok{ifelse}\NormalTok{(df}\SpecialCharTok{$}\NormalTok{OUTCOME }\SpecialCharTok{==} \StringTok{"dism\_other"}\NormalTok{, }\DecValTok{1}\NormalTok{, }\DecValTok{0}\NormalTok{)}
\end{Highlighting}
\end{Shaded}

Let's check if we did it correctly.

\begin{Shaded}
\begin{Highlighting}[]
\NormalTok{df }\SpecialCharTok{\%\textgreater{}\%} \FunctionTok{select}\NormalTok{(OUTCOME, dism\_other) }\SpecialCharTok{\%\textgreater{}\%} \FunctionTok{head}\NormalTok{()}
\end{Highlighting}
\end{Shaded}

\begin{verbatim}
## # A tibble: 6 x 2
##   OUTCOME               dism_other
##   <fct>                      <dbl>
## 1 jud_default_plaintiff          0
## 2 dism_other                     1
## 3 dism_pros                      0
## 4 dism_other                     1
## 5 dism_vol                       0
## 6 dism_vol                       0
\end{verbatim}

But there is a total of 34 categories in the \texttt{OUTCOME} variable,
it would be a waste of time to generate all 34 dummies if we have to
retype the above one by one for each category. Luckily, R is here to
help us! We can run a \texttt{forloop} function to ask R to do the same
thing for all the 34 categories.

\begin{Shaded}
\begin{Highlighting}[]
\ControlFlowTok{for}\NormalTok{(v }\ControlFlowTok{in} \FunctionTok{levels}\NormalTok{(df}\SpecialCharTok{$}\NormalTok{OUTCOME))\{}
  \FunctionTok{print}\NormalTok{(v)}
\NormalTok{  df[[v]] }\OtherTok{\textless{}{-}} \FunctionTok{ifelse}\NormalTok{(df}\SpecialCharTok{$}\NormalTok{OUTCOME }\SpecialCharTok{==}\NormalTok{ v, }\DecValTok{1}\NormalTok{, }\DecValTok{0}\NormalTok{) }\DocumentationTok{\#\# df[[v]] is a way to call out/create a variable}
\NormalTok{\} }\DocumentationTok{\#\# This is called "forloop."}
\end{Highlighting}
\end{Shaded}

\begin{verbatim}
## [1] "dism_juris"
## [1] "dism_other"
## [1] "dism_pros"
## [1] "dism_vol"
## [1] "jud_bench_both"
## [1] "jud_bench_defendant"
## [1] "jud_bench_NA"
## [1] "jud_bench_plaintiff"
## [1] "jud_consent_both"
## [1] "jud_consent_defendant"
## [1] "jud_consent_NA"
## [1] "jud_consent_plaintiff"
## [1] "jud_default_both"
## [1] "jud_default_defendant"
## [1] "jud_default_NA"
## [1] "jud_default_plaintiff"
## [1] "jud_directed_defendant"
## [1] "jud_directed_NA"
## [1] "jud_directed_plaintiff"
## [1] "jud_jury_both"
## [1] "jud_jury_defendant"
## [1] "jud_jury_NA"
## [1] "jud_jury_plaintiff"
## [1] "jud_misc_both"
## [1] "jud_misc_defendant"
## [1] "jud_misc_NA"
## [1] "jud_misc_plaintiff"
## [1] "jud_motion_both"
## [1] "jud_motion_defendant"
## [1] "jud_motion_NA"
## [1] "jud_motion_plaintiff"
## [1] "remand"
## [1] "settlement"
## [1] "transfer"
\end{verbatim}

\begin{Shaded}
\begin{Highlighting}[]
\FunctionTok{ncol}\NormalTok{(df)}
\end{Highlighting}
\end{Shaded}

\begin{verbatim}
## [1] 88
\end{verbatim}

You can find that the number of columns increased to 88 from 54 after
creating 34 dummies.

\hypertarget{g.-the-rank-of-a-matrix}{%
\subparagraph{g. The rank of a matrix}\label{g.-the-rank-of-a-matrix}}

The rank of a dataset is the number of ``linearly independent''
variables.

\begin{Shaded}
\begin{Highlighting}[]
\FunctionTok{library}\NormalTok{(Matrix) }\DocumentationTok{\#\# This package has a function for calculating the rank of a matrix}

\CommentTok{\# Let\textquotesingle{}s create a fake dataset}
\NormalTok{df2 }\OtherTok{\textless{}{-}} \FunctionTok{data.frame}\NormalTok{(}\AttributeTok{var1 =} \FunctionTok{seq}\NormalTok{(}\DecValTok{1}\NormalTok{, }\DecValTok{100}\NormalTok{, }\DecValTok{3}\NormalTok{)) }\DocumentationTok{\#\# Creates 34 rows}

\NormalTok{df2}\SpecialCharTok{$}\NormalTok{var2 }\OtherTok{\textless{}{-}}\NormalTok{ (df2}\SpecialCharTok{$}\NormalTok{var1}\SpecialCharTok{/}\NormalTok{pi) }\DocumentationTok{\#\# create a new variable called var2 in df2}
\NormalTok{pi }\DocumentationTok{\#\# pi is an inbuilt R constant whose value is 3.141593}
\end{Highlighting}
\end{Shaded}

\begin{verbatim}
## [1] 3.141593
\end{verbatim}

\begin{Shaded}
\begin{Highlighting}[]
\NormalTok{df2}\SpecialCharTok{$}\NormalTok{var3 }\OtherTok{\textless{}{-}} \DecValTok{6} 
\NormalTok{df2}
\end{Highlighting}
\end{Shaded}

\begin{verbatim}
##    var1       var2 var3
## 1     1  0.3183099    6
## 2     4  1.2732395    6
## 3     7  2.2281692    6
## 4    10  3.1830989    6
## 5    13  4.1380285    6
## 6    16  5.0929582    6
## 7    19  6.0478878    6
## 8    22  7.0028175    6
## 9    25  7.9577472    6
## 10   28  8.9126768    6
## 11   31  9.8676065    6
## 12   34 10.8225361    6
## 13   37 11.7774658    6
## 14   40 12.7323954    6
## 15   43 13.6873251    6
## 16   46 14.6422548    6
## 17   49 15.5971844    6
## 18   52 16.5521141    6
## 19   55 17.5070437    6
## 20   58 18.4619734    6
## 21   61 19.4169031    6
## 22   64 20.3718327    6
## 23   67 21.3267624    6
## 24   70 22.2816920    6
## 25   73 23.2366217    6
## 26   76 24.1915513    6
## 27   79 25.1464810    6
## 28   82 26.1014107    6
## 29   85 27.0563403    6
## 30   88 28.0112700    6
## 31   91 28.9661996    6
## 32   94 29.9211293    6
## 33   97 30.8760590    6
## 34  100 31.8309886    6
\end{verbatim}

What is the rank of this matrix (the fake dataset we created)?

\begin{Shaded}
\begin{Highlighting}[]
\FunctionTok{rankMatrix}\NormalTok{(}\FunctionTok{as.matrix}\NormalTok{(df2)) }\DocumentationTok{\#\# The input must be a numeric matrix, which means we have to covert df2 from a dataframe to a matrix}
\end{Highlighting}
\end{Shaded}

\begin{verbatim}
## [1] 2
## attr(,"method")
## [1] "tolNorm2"
## attr(,"useGrad")
## [1] FALSE
## attr(,"tol")
## [1] 7.549517e-15
\end{verbatim}

There are 3 variables in df2, but there are only 2 linearly independent
variables since the generation of var2 is based on var1.

\end{document}
